Khóa luận này đề xuất kiến trúc Graph-based Retrieval-Augmented Generation (GraphRAG) cho bài toán khám phá và truy vấn tri thức trong miền nghệ thuật Chèo - một loại hình sân khấu truyền thống của Việt Nam có cấu trúc quan hệ phức tạp giữa các vở diễn, nhân vật, diễn viên và các cảnh diễn. Ba đóng góp chính của nghiên cứu gồm: (i) xây dựng đồ thị tri thức Chèo cho bảy vở chèo cổ với 38 nhân vật, 41 diễn viên và 15 trích đoạn, mở rộng từ ontology gốc với 14 luật suy diễn quan hệ gián tiếp và ba mức chỉ mục phục vụ các truy vấn phức tạp; (ii) đề xuất kiến trúc GraphRAG với hai phân hệ G-Retrieval (truy vấn đa mức độ trên đồ thị) và G-Generation (ba chiến lược tạo sinh Pre/Mid/Post) tích hợp cơ chế chống ảo giác qua Danh mục thực thể; và (iii) thiết kế bộ câu hỏi đánh giá CheoBench gồm 100 câu hỏi phân tầng theo ba mức Local, Community, Global cùng hệ độ đo chín thành phần thuộc ba khối Information Retrieval (IR), Context và Generation.

Kết quả thực nghiệm trên CheoBench cho thấy GraphRAG đạt điểm tổng hợp $0,850$, vượt RAG truyền thống ($0,613$) chủ yếu ở khối Context (Context Recall $0,964$ so với $0,418$) và Generation (Faithfulness $0,986$ so với $0,734$), trong khi tương đương ở khối IR thuần - lợi ích của cấu trúc đồ thị nằm ở giai đoạn tổ chức ngữ cảnh, không phải ở pha truy xuất. Khảo sát $103$ người dùng độc lập cho thấy kết quả mang tính định lượng : $89,6\%$ người dùng chọn GraphRAG là hệ thống tốt nhất ở khảo sát chấm điểm đa tiêu chí, khẳng định tính khả thi của việc tích hợp đồ thị tri thức vào kiến trúc Retrieval-Augmented Generation (RAG) cho miền tri thức có cấu trúc quan hệ phức tạp.

\vspace{0.5em}
\noindent\textbf{Từ khoá:} GraphRAG, đồ thị tri thức, Retrieval-Augmented Generation, nghệ thuật Chèo, hỏi đáp tri thức.
